% SCP10, Definition 6.13 and Lemma 6.14, paper_v3.tex:2174--2214;
% detection comparison: Theorem 6.15, lines 2216--2267.
% Exact scope: finite generated dual homotopies on labelled periodic lattices.
% This fragment does not mark the unrestricted source lemma as formalized.

\begin{definition}[Oriented dual paths and signed crossings]%
    \label{def:peps_dual_flux_path}
    \lean{TNLean.PEPS.TorusDualStep,
        TNLean.PEPS.torusDualQuiver,
        TNLean.PEPS.TorusDualPath,
        TNLean.PEPS.TorusDualStep.reverse,
        TNLean.PEPS.TorusDualStep.reverse_reverse,
        TNLean.PEPS.torusDualInvolutiveReverse,
        TNLean.PEPS.torusPointMass,
        TNLean.PEPS.TorusDualStep.crossings,
        TNLean.PEPS.TorusDualStep.crossings_reverse,
        TNLean.PEPS.torusDualPathCrossings,
        TNLean.PEPS.TorusDualPath.crossings_nil,
        TNLean.PEPS.TorusDualPath.crossings_cons,
        TNLean.PEPS.TorusDualPath.crossings_comp,
        TNLean.PEPS.TorusDualPath.crossings_reverse}
    \leanok
    Let $X=\mathbb Z_w\times\mathbb Z_h$, where $w,h>0$.
    A plaquette is labelled by its lower-left primal vertex.
    A finite dual path is a sequence of east, west, north and south
    unit steps, retaining the step labels at periodic seams and when
    parallel bonds have the same endpoints. Write $n(p)=(n_h(p),n_v(p))$
    for its integer crossing counts on rightward and downward native bonds.
    With $\delta_s(v)=\mathbf1_{\{v=s\}}$, the positive steps contribute
    \begin{align}
        n\bigl((x,y)\to(x+1,y)\bigr) & =(0,\delta_{(x+1,y)}), &
        n\bigl((x,y)\to(x,y+1)\bigr) & =(\delta_{(x,y+1)},0).
        \label{eq:peps_dual_step_counts}
    \end{align}
    Opposite steps negate these counts; concatenation adds them.
    Thus $n(p^{-1})=-n(p)$, with multiplicities retained.
    In the orientation of \cite[Definition~6.13]{Schuch2010PEPS},
    an eastward dual step inserts $U_g$ on a downward bond and a
    southward step inserts $U_{g^{-1}}$ on a rightward bond.
\end{definition}

\begin{definition}[Finite deformations with prescribed swept vertices]%
    \label{def:peps_dual_flux_homotopy}
    \lean{TNLean.PEPS.torusFluxGradient,
        TNLean.PEPS.torusFluxGradient_zero,
        TNLean.PEPS.torusFluxGradient_add,
        TNLean.PEPS.torusFluxGradient_neg,
        TNLean.PEPS.torusDualEastNorth,
        TNLean.PEPS.torusDualNorthEast,
        TNLean.PEPS.TorusDualHomotopy,
        TNLean.PEPS.torusDualSquareLoop}
    \leanok
    \uses{def:peps_dual_flux_path}
    For $R\subseteq X$, let $p\sim_R q$ be the relation generated by
    removal of an adjacent step and its reverse and by the square move
    \begin{align}
        (x,y)\to(x+1,y)\to(x+1,y+1)
        \quad\longleftrightarrow\quad
        (x,y)\to(x,y+1)\to(x+1,y+1),
        \label{eq:peps_dual_square_move}
    \end{align}
    allowed when $(x+1,y+1)\in R$. The relation is closed under
    reflexivity, symmetry, transitivity, reversal and concatenation.
    The endpoints remain fixed. For $f:X\to\mathbb Z$, define
    \begin{align}
        (\nabla f)_h(v) & =f(v+(1,0))-f(v), &
        (\nabla f)_v(v) & =f(v)-f(v+(0,1)).
        \label{eq:peps_dual_gradient}
    \end{align}
    This describes the elementary deformations used in
    \cite[proof of Lemma~6.14]{Schuch2010PEPS}.
\end{definition}

% Formula: n(NE)-n(EN)=grad(delta_s), s=(x+1,y+1).
% Source: SCP10 figs5/def-flux-pair.pdf and eq:2d:move-strings.
% Ink-to-index: centre A_s; solid horizontal arrows point right and vertical
% arrows down. EN crosses the bottom and right native bonds; NE crosses the
% left and top bonds. Each crossed bond contains the actual matrix U_g.
% The dotted path is a dual geometric route, not an additional tensor index.
% Contracted legs: each U_g has its inner virtual index summed with A_s.
% Open boundary: four virtual indices and one physical index per panel.
% Both dual routes have the same two geometric endpoints a,b.
\begin{tenkzequation}
    \tenkzkernel
    \begin{tenkz}[size=m, rows={wire,wire,wire,wire,wire,wire,wire}, cols=7, bonds=none]
        \tn[at=(4,4), name=fluxA, skin=box,
            ports={0:virtual,90:virtual,180:virtual,270:virtual,45:physical:$\sigma_s$}]{A_s}
        \tn[at=(4,6), name=fluxE, skin=box, ports={0:virtual,180:virtual}]{U_g}
        \tn[at=(6,4), name=fluxS, skin=box, ports={90:virtual,270:virtual}]{U_g}
        \tnwire[dir=from]{fluxA.180}{open w}
        \tnwire[dir=from]{fluxA.90}{open n}
        \tnwire[dir=to, name=fluxEin]{fluxA.0}{fluxE.180}
        \tnwire[dir=to, name=fluxEout]{fluxE.0}{open e}
        \tnwire[dir=to, name=fluxSin]{fluxA.270}{fluxS.90}
        \tnwire[dir=to, name=fluxSout]{fluxS.270}{open s}
        \tnwire[kind=string, stroke=dotted, route=orth, dir=to,
            name=fluxEN, via={(6,6)},
            cross={under at crossing of fluxEN and fluxEin,
                    under at crossing of fluxEN and fluxEout,
                    under at crossing of fluxEN and fluxSin,
                    under at crossing of fluxEN and fluxSout}]{(6,2)}{(2,6)}
        \tnmark[form=label, label pos=225]{(6,2)}{$a$}
        \tnmark[form=label, label pos=45]{(2,6)}{$b$}
    \end{tenkz}
    $\qquad\longleftrightarrow\qquad$
    \begin{tenkz}[size=m, rows={wire,wire,wire,wire,wire,wire,wire}, cols=7, bonds=none]
        \tn[at=(4,4), name=fluxA, skin=box,
            ports={0:virtual,90:virtual,180:virtual,270:virtual,45:physical:$\sigma_s$}]{A_s}
        \tn[at=(4,2), name=fluxW, skin=box, ports={0:virtual,180:virtual}]{U_g}
        \tn[at=(2,4), name=fluxN, skin=box, ports={90:virtual,270:virtual}]{U_g}
        \tnwire[dir=from, name=fluxWin]{fluxA.180}{fluxW.0}
        \tnwire[dir=from, name=fluxWout]{fluxW.180}{open w}
        \tnwire[dir=from, name=fluxNin]{fluxA.90}{fluxN.270}
        \tnwire[dir=from, name=fluxNout]{fluxN.90}{open n}
        \tnwire[dir=to]{fluxA.0}{open e}
        \tnwire[dir=to]{fluxA.270}{open s}
        \tnwire[kind=string, stroke=dotted, route=orth, dir=to,
            name=fluxNE, via={(2,2)},
            cross={under at crossing of fluxNE and fluxWin,
                    under at crossing of fluxNE and fluxWout,
                    under at crossing of fluxNE and fluxNin,
                    under at crossing of fluxNE and fluxNout}]{(6,2)}{(2,6)}
        \tnmark[form=label, label pos=225]{(6,2)}{$a$}
        \tnmark[form=label, label pos=45]{(2,6)}{$b$}
    \end{tenkz}
\end{tenkzequation}
The two routes in (\ref{eq:peps_dual_square_move}) cross opposite pairs
of bonds incident to $s=(x+1,y+1)$.

\begin{theorem}[A potential derived from the deformation]%
    \label{thm:peps_dual_flux_potential}
    \lean{TNLean.PEPS.torusDualSquare_crossings,
        TNLean.PEPS.TorusDualHomotopy.exists_potential,
        TNLean.PEPS.TorusDualHomotopy.cancel_reverse,
        TNLean.PEPS.TorusDualHomotopy.squareLoop}
    \leanok
    \uses{def:peps_dual_flux_homotopy}
    If $p\sim_R q$, there is an integer function $f$ vanishing outside
    $R$ such that $n(q)=n(p)+\nabla f$. A square move has
    $f=\delta_{(x+1,y+1)}$. Retracing any finite path contracts by
    backtrack removal, and the four-sided dual square loop contracts
    by one square move followed by backtrack removal.
\end{theorem}
\begin{proof}\leanok
    Substitution of (\ref{eq:peps_dual_step_counts}) into
    (\ref{eq:peps_dual_gradient}) gives the square identity.
    Induct on the generating moves: backtracks have potential zero,
    symmetry and reversal negate a potential, and composition and
    transitivity add potentials. Each operation preserves vanishing
    outside $R$.
\end{proof}

\begin{definition}[Lifted displacement]%
    \label{def:peps_dual_flux_displacement}
    \lean{TNLean.PEPS.TorusDualStep.displacement,
        TNLean.PEPS.TorusDualStep.displacement_reverse,
        TNLean.PEPS.torusDualPathDisplacement,
        TNLean.PEPS.torusDualPathDisplacement_nil,
        TNLean.PEPS.torusDualPathDisplacement_cons,
        TNLean.PEPS.torusDualPathDisplacement_comp,
        TNLean.PEPS.torusDualPathDisplacement_reverse,
        TNLean.PEPS.TorusDualStep.displacement_cast,
        TNLean.PEPS.torusDualPathDisplacement_cast}
    \leanok
    \uses{def:peps_dual_flux_path}
    Sum $(1,0),(-1,0),(0,1),(0,-1)$ over the east, west, north
    and south steps of $p$ to obtain $d(p)\in\mathbb Z^2$.
    Concatenation adds displacements, and reversal negates them.
    For $p:a\to b$, reducing the two coordinates of $d(p)$ modulo
    $w$ and $h$ gives $b-a\in X$.
\end{definition}
\begin{theorem}[Winding is preserved]%
    \label{thm:peps_dual_flux_displacement}
    \lean{TNLean.PEPS.TorusDualHomotopy.displacement_eq}
    \leanok
    \uses{def:peps_dual_flux_homotopy,def:peps_dual_flux_displacement}
    If $p\sim_R q$, then $d(p)=d(q)$. Thus equal torus endpoints
    alone do not imply that two paths are related by these moves.
\end{theorem}
\begin{proof}\leanok
    A backtrack has zero displacement, and both sides of a square
    move have displacement $(1,1)$. Induction over the relation uses
    additivity and the sign change under reversal.
\end{proof}

\begin{definition}[Native flux matrices and unchanged exterior data]%
    \label{def:peps_dual_flux_matrices}
    \lean{TNLean.PEPS.torusFluxPowerLabels,
        TNLean.PEPS.torusDualFluxLabels,
        TNLean.PEPS.torusDualFluxState,
        TNLean.PEPS.torusFluxMatricesWithExterior}
    \leanok
    \uses{def:peps_dual_flux_path}
    Let $G$ be a group, $g\in G$, and $U:G\to\operatorname{GL}(V)$
    a representation on a finite-dimensional virtual space with a fixed
    basis. Assign $u_h^p(v)=g^{n_h(p)(v)}$ and
    $u_v^p(v)=g^{n_v(p)(v)}$ to the two native bond families, and put
    the actual matrices $U_{u_h^p(v)}$, $U_{u_v^p(v)}$ on those bonds.
    For a swept set $R$, one may instead retain these matrices only
    on bonds incident to $R$ and prescribe arbitrary fixed matrices
    $E_e$ on every bond whose two endpoints lie outside $R$.
    Denote the actual closed PEPS coefficient, with site tensors
    $(A_v)$ and physical configuration $\sigma$, by $\Psi_p^E(\sigma)$.
\end{definition}

\begin{theorem}[Absorption of a supported vertex gauge]%
    \label{thm:peps_dual_flux_gauge_absorption}
    \lean{TNLean.PEPS.torusBondNetwork_groupGauge_of_vecMul_eq,
        TNLean.PEPS.torusBondNetwork_groupGauge_of_invariant_on,
        TNLean.PEPS.torusBondNetwork_groupGauge,
        TNLean.PEPS.torusBondNetwork_torusBondGauge,
        TNLean.PEPS.torusBondNetwork_torusBondGauge_of_invariant_on,
        TNLean.PEPS.torusBondGauge_fst_eq_of_eq_one,
        TNLean.PEPS.torusBondGauge_snd_eq_of_eq_one,
        TNLean.PEPS.torusBondGauge_fst_eq_of_not_mem,
        TNLean.PEPS.torusBondGauge_snd_eq_of_not_mem}
    \leanok
    For arbitrary native bond matrices $(O_h,O_v)$ and $k:X\to G$,
    make the substitutions
    \begin{align}
        O'_h(v) & =U_{k(v+(1,0))}O_h(v)U_{k(v)^{-1}}, &
        O'_v(v) & =U_{k(v)}O_v(v)U_{k(v+(0,1))^{-1}}.
        \label{eq:peps_dual_matrix_gauge}
    \end{align}
    The closed contraction is unchanged if each site tensor absorbs
    the induced four-leg action of its $k(v)$. It suffices to assume
    virtual $G$-invariance only at sites where $k(v)\ne1$.
    If $k=1$ outside $R$, all bonds with both endpoints outside $R$
    and all exterior site tensors remain unchanged.
\end{theorem}
\begin{proof}\leanok
    Expand each matrix product in (\ref{eq:peps_dual_matrix_gauge})
    inside the finite bond sum. Group the four adjacent factors with
    each site tensor and absorb them using its virtual invariance.
    At a site with $k(v)=1$ all four factors are identities.
\end{proof}

\begin{theorem}[Physical invariance under finite dual-path deformation]%
    \label{thm:peps_dual_flux_deformation}
    \lean{TNLean.PEPS.torusFluxPowerLabels_add_gradient,
        TNLean.PEPS.TorusDualHomotopy.exists_fluxGauge,
        TNLean.PEPS.TorusDualHomotopy.torusBondNetwork_eq,
        TNLean.PEPS.TorusDualHomotopy.torusDualFluxState_eq,
        TNLean.PEPS.TorusDualHomotopy.torusBondNetwork_eq_with_exterior}
    \leanok
    \uses{def:peps_dual_flux_homotopy,def:peps_dual_flux_matrices}
    Suppose $p\sim_R q$ and every $A_v$ with $v\in R$ is virtually
    $G$-invariant. Then $\Psi_p^E(\sigma)=\Psi_q^E(\sigma)$ for every
    physical configuration and every unchanged exterior assignment $E$.
    The site tensors and their physical spaces may vary with $v$;
    exterior tensors are arbitrary. This holds for all positive periods,
    including one and two, and for nonabelian $G$.
    It is the finite generated-deformation form of
    \cite[Lemma~6.14]{Schuch2010PEPS}.
\end{theorem}
\begin{proof}\leanok
    \uses{thm:peps_dual_flux_potential,thm:peps_dual_flux_gauge_absorption}
    Take the derived potential $f$ and set $k(v)=g^{f(v)}$.
    The exponent identity gives
    \begin{align}
        u_h^q(v) & =k(v+(1,0))u_h^p(v)k(v)^{-1}, &
        u_v^q(v) & =k(v)u_v^p(v)k(v+(0,1))^{-1}.
        \label{eq:peps_dual_label_gauge}
    \end{align}
    Only powers of the single element $g$ are combined, so group
    commutativity is unnecessary. Since $k=1$ outside $R$, arbitrary
    exterior matrices satisfy (\ref{eq:peps_dual_matrix_gauge}) as well.
    Apply Theorem~\ref{thm:peps_dual_flux_gauge_absorption} to the actual
    closed contraction.
\end{proof}

\begin{definition}[Clockwise native holonomy and signed curl]%
    \label{def:peps_dual_flux_curl}
    \lean{TNLean.PEPS.torusFluxCurl,
        TNLean.PEPS.torusFluxCurl_zero,
        TNLean.PEPS.torusFluxCurl_add,
        TNLean.PEPS.torusBondPlaquetteHolonomy,
        TNLean.PEPS.torusBondPlaquetteHolonomy_zpow}
    \leanok
    \uses{def:peps_dual_flux_path}
    For native group labels $u=(u_h,u_v)$, define the clockwise
    holonomy at $z=(x,y)$ by
    \begin{align}
        H_u(z)                    & =u_h(z)^{-1}u_v(x+1,y)u_h(x,y+1)u_v(z)^{-1},
        \label{eq:peps_dual_native_holonomy}                                     \\
        (\operatorname{curl}n)(z) & =n_h(x,y+1)+n_v(x+1,y)-n_h(z)-n_v(z).
        \label{eq:peps_dual_native_curl}
    \end{align}
    Curl is additive. For the labels $u=(g^{n_h},g^{n_v})$,
    $H_u(z)=g^{(\operatorname{curl}n)(z)}$.
\end{definition}

\begin{theorem}[Flux is supported at the two path endpoints]%
    \label{thm:peps_dual_flux_endpoints}
    \lean{TNLean.PEPS.TorusDualStep.fluxCurl,
        TNLean.PEPS.TorusDualPath.fluxCurl,
        TNLean.PEPS.torusBondPlaquetteHolonomy_torusDualFluxLabels,
        TNLean.PEPS.torusBondPlaquetteHolonomy_torusDualFluxLabels_source,
        TNLean.PEPS.torusBondPlaquetteHolonomy_torusDualFluxLabels_target,
        TNLean.PEPS.torusBondPlaquetteHolonomy_torusDualFluxLabels_of_ne,
        TNLean.PEPS.torusBondPlaquetteHolonomy_torusDualFluxLabels_loop}
    \leanok
    \uses{def:peps_dual_flux_curl,def:peps_dual_flux_matrices}
    Every finite dual path $p:a\to b$ satisfies
    \begin{align}
        \operatorname{curl}n(p) & =\delta_a-\delta_b,           &
        H_{u^p}(z)              & =g^{\delta_a(z)-\delta_b(z)}.
        \label{eq:peps_dual_endpoint_flux}
    \end{align}
    If $a\ne b$, the clockwise labels are $g$ at $a$, $g^{-1}$ at
    $b$, and identity elsewhere, as in
    \cite[Definition~6.13]{Schuch2010PEPS}. A closed dual path has
    identity plaquette holonomy everywhere, including when it winds
    around the torus.
\end{theorem}
\begin{proof}\leanok
    A single step contributes $\delta_a-\delta_b$ by
    (\ref{eq:peps_dual_step_counts}). Add along the path; each intermediate
    endpoint occurs once with each sign. Exponentiation gives
    (\ref{eq:peps_dual_endpoint_flux}).
\end{proof}

\begin{theorem}[Clockwise and counterclockwise flux conventions]%
    \label{thm:peps_dual_flux_orientation}
    \lean{TNLean.PEPS.torusNativeRightTransport_torusGraphRegularLabels,
        TNLean.PEPS.torusNativeUpTransport_torusGraphRegularLabels,
        TNLean.PEPS.regularWalkHolonomy_torusGraphRegularLabels_reverse,
        TNLean.PEPS.regularWalkHolonomy_torusGraphRegularLabels,
        TNLean.PEPS.regularWalkHolonomy_torusGraphRegularLabels_zpow_reverse,
        TNLean.PEPS.regularWalkHolonomy_torusGraphRegularLabels_zpow,
        TNLean.PEPS.regularWalkHolonomy_torusDualFluxLabels_reverse,
        TNLean.PEPS.regularWalkHolonomy_torusDualFluxLabels}
    \leanok
    \uses{def:peps_dual_flux_curl,thm:peps_dual_flux_endpoints}
    Suppose $w,h\geq3$. Passing from native bonds to the ordered edges
    of the torus graph gives rightward transport $u_h(v)$ and upward
    transport $u_v(v)^{-1}$, including on both seams. Hence the
    counterclockwise plaquette walk has holonomy $H_u(z)^{-1}$ and
    its reverse has holonomy $H_u(z)$. For a dual path, the
    counterclockwise exponent is $\delta_b(z)-\delta_a(z)$.
\end{theorem}
\begin{proof}\leanok
    Compare each ordered edge with its native direction. When the
    directions disagree, taking the inverse in the directed transport
    cancels the inverse in the edge label. The four-step products
    then give (\ref{eq:peps_dual_native_holonomy}) and its inverse.
\end{proof}

\begin{definition}[The physical cut matrix of native regular insertions]%
    \label{def:peps_dual_flux_cut}
    \lean{TNLean.PEPS.torusBondRegularPhysicalCutMatrix,
        TNLean.PEPS.torusBondNetwork_leftRegular_eq_stateCoeff,
        TNLean.PEPS.torusBondRegularPhysicalCutMatrix_eq_regularPhysicalCutMatrix}
    \leanok
    \uses{def:peps_dual_flux_matrices}
    For finite $G$, $w,h\geq3$, and a homogeneous tensor $A$, use
    the left-regular matrices $U_g=L_g$. For a physical region $S$,
    let $M_{u,S}(\sigma_S,\sigma_{S^c})$ be the actual native
    contraction with these physical indices fixed. It agrees entrywise
    with the cut matrix of the ordered-edge regular tensor network.
    Write $S_z$ for the four original physical sites of plaquette $z$.
\end{definition}
\begin{proof}\leanok
    Reindex the native bond-end sum by the ordered graph edges.
    The regular matrix on each edge imposes precisely the translated
    equality of its two labels, leaving the regular twisted graph
    contraction. Apply this coefficient identity to the assembled
    physical configuration on $S$ and its complement.
\end{proof}

\begin{theorem}[Nonvanishing of native regular flux strings]%
    \label{thm:peps_dual_flux_nonzero}
    \lean{TNLean.PEPS.IsGIsometric.torusDualFlux_state_ne_zero,
        TNLean.PEPS.IsGIsometric.torusDualFlux_cutMatrix_ne_zero}
    \leanok
    \uses{def:peps_dual_flux_cut}
    If $A$ is regular $G$-isometric, every actual native dual-string
    state $\Psi_p$ is nonzero. Its coefficient matrix $M_{u^p,S}$
    across every physical cut $S$ is also nonzero.
\end{theorem}
\begin{proof}\leanok
    \uses{def:peps_dual_flux_cut,thm:peps_regular_twisted_state_nonzero}
    The coefficient identification gives the regular twisted graph
    state, which is nonzero for every edge assignment. If a cut matrix
    vanished, evaluate it at the restrictions of any global physical
    configuration to recover a zero coefficient at that configuration.
    The whole state would then vanish, a contradiction.
\end{proof}

\begin{theorem}[Four-spin detection of native flux labels]%
    \label{thm:peps_dual_flux_detection}
    \lean{TNLean.PEPS.IsGIsometric.exists_torusNative_plaquette_cutMeasurement,
        TNLean.PEPS.IsGIsometric.exists_torusNativePower_plaquette_cutMeasurement,
        TNLean.PEPS.IsGIsometric.exists_torusDualFlux_counterclockwise_cutMeasurement,
        TNLean.PEPS.conjClassesInvEquiv,
        TNLean.PEPS.conjClassesInvEquiv_mk,
        TNLean.PEPS.IsGIsometric.exists_torusDualFlux_clockwise_cutMeasurement}
    \leanok
    \uses{def:peps_dual_flux_cut,thm:peps_dual_flux_orientation,
        thm:peps_dual_flux_nonzero}
    Let $A$ be regular $G$-isometric. On $S_z$ there is one complete
    projective measurement $(Q_C^{\rm ccw})_C$, with an additional
    outcome $Q_\perp^{\rm ccw}$, such that for every native assignment $u$
    \begin{align}
        Q_C^{\rm ccw}M_{u,S_z}                   & =
        \mathbf1_{\{C=[H_u(z)^{-1}]\}}M_{u,S_z}, &
        Q_\perp^{\rm ccw}M_{u,S_z}               & =0.
        \label{eq:peps_dual_ccw_measurement}
    \end{align}
    All projectors are Hermitian and positive, distinct outcomes are
    orthogonal, and their sum is the identity. Inversion permutes
    conjugacy classes by $\iota([s])=[s^{-1}]$. Relabelling
    $Q_C^{\rm cw}=Q_{\iota(C)}^{\rm ccw}$ and leaving the extra outcome
    unchanged gives, simultaneously for every $g$ and every finite
    dual path $p:a\to b$,
    \begin{align}
        Q_C^{\rm cw}M_{u^p,S_z}                                    & =
        \mathbf1_{\{C=[g^{\delta_a(z)-\delta_b(z)}]\}}M_{u^p,S_z}, &
        Q_\perp^{\rm cw}M_{u^p,S_z}                                & =0.
        \label{eq:peps_dual_cw_measurement}
    \end{align}
    Thus the source convention detects $[g]$ and $[g^{-1}]$ at the
    distinct endpoints of every such state. The two classes
    need not coincide. These are the native regular torus instances of
    \cite[Theorem~6.15]{Schuch2010PEPS}.
\end{theorem}
\begin{proof}\leanok
    \uses{thm:peps_torus_plaquette_holonomy_measurement,
        def:peps_dual_flux_cut,thm:peps_dual_flux_orientation,
        thm:peps_dual_flux_endpoints}
    Transfer the fixed plaquette projectors through the entrywise cut
    identity and substitute the inverse holonomy formula. Inversion
    preserves conjugacy and is involutive, so relabelling preserves
    positivity, orthogonality and completeness. Finally substitute
    (\ref{eq:peps_dual_endpoint_flux}) into
    (\ref{eq:peps_dual_ccw_measurement}).
\end{proof}

\begin{theorem}[Local parent constraints at identity-flux plaquettes]%
    \label{thm:peps_dual_flux_local_parent}
    \lean{TNLean.PEPS.exists_torusBondGauge_plaquette_internal_eq_one,
        TNLean.PEPS.torusBondNetwork_slice_mem_plaquetteGroundSpace_of_holonomy_eq_one,
        TNLean.PEPS.torusParent_annihilates_torusBondNetwork_of_holonomy_eq_one,
        TNLean.PEPS.canonicalTorusParent_annihilates_torusBondNetwork_of_holonomy_eq_one,
        TNLean.PEPS.torusParent_annihilates_torusBondNetwork_zpow_of_curl_eq_zero}
    \leanok
    \uses{def:peps_dual_flux_curl,def:peps_dual_flux_matrices}
    Let $w,h\geq3$, and let $A$ be a homogeneous virtually
    $G$-invariant tensor for an arbitrary finite-dimensional
    representation $U$. If $H_u(z)=1$, every complementary physical
    slice of the inserted state lies in the original plaquette PEPS
    range $\mathcal G_{S_z}$. Consequently, for every positive local
    interaction $P_z$ with kernel $\mathcal G_{S_z}$,
    \begin{align}
        (P_z\otimes I_{S_z^c})\Psi_u & =0.
        \label{eq:peps_dual_parent_constraint}
    \end{align}
    This includes the canonical orthogonal projector onto
    $\mathcal G_{S_z}^{\perp}$. In particular, zero signed curl of
    native powers implies (\ref{eq:peps_dual_parent_constraint}).
\end{theorem}
\begin{proof}\leanok
    \uses{thm:peps_dual_flux_gauge_absorption}
    With $z=(x,y)$, put $k(z)=1$,
    $k(x+1,y)=u_h(z)^{-1}$, $k(x,y+1)=u_v(z)$ and
    $k(x+1,y+1)=u_v(z)u_h(x,y+1)^{-1}$, and set $k=1$ elsewhere.
    The condition $H_u(z)=1$ makes all four internal transformed
    bonds identities. Gauge absorption preserves the actual state.
    The remaining crossing and exterior contractions provide a
    boundary vector, so each complementary slice belongs to
    $\mathcal G_{S_z}$. Apply the defining kernel property of $P_z$.
\end{proof}


\begin{theorem}[Parent interactions away from dual-string endpoints]%
    \label{thm:peps_dual_flux_endpoint_parent}
    \lean{TNLean.PEPS.torusDualFluxState_parent_annihilates_of_ne,
        TNLean.PEPS.torusDualFluxState_canonicalParent_annihilates_of_ne,
        TNLean.PEPS.torusDualFluxState_canonicalParent_annihilates_loop}
    \leanok
    \uses{thm:peps_dual_flux_endpoints,thm:peps_dual_flux_local_parent}
    Under the hypotheses of Theorem~\ref{thm:peps_dual_flux_local_parent},
    let $p:a\to b$ be any finite dual path. Every positive parent
    interaction at a plaquette $z\notin\{a,b\}$ annihilates the actual
    inserted state $\Psi_p$, including the canonical parent projector.
    If $a=b$, every canonical plaquette parent term annihilates $\Psi_p$.
\end{theorem}
\begin{proof}\leanok
    \uses{thm:peps_dual_flux_endpoints,thm:peps_dual_flux_local_parent}
    Formula~(\ref{eq:peps_dual_endpoint_flux}) supplies identity holonomy
    at each indicated plaquette. Apply
    Theorem~\ref{thm:peps_dual_flux_local_parent}.
\end{proof}

% SCP10 arXiv:1001.3807v3, Lemma 6.14, paper_v3.tex:2199--2214.
% Candidate finite-rectangle specialization. Lean compilation is blocked by
% unavailable pinned Mathlib cache; intentionally no leanok or lean tags yet.
\begin{theorem}[Deformation inside an embedded rectangular dual patch]%
    \label{thm:peps_dual_rectangle}
    \notready
    \uses{def:peps_dual_flux_homotopy}
    Fix a dual plaquette $(x_0,y_0)$ and integers $0\leq W<w$,
    $0\leq H<h$. Embed the closed rectangle
    $D=\{0,\ldots,W\}\times\{0,\ldots,H\}$ by
    $\iota(i,j)=(x_0+i,y_0+j)$ modulo $(w,h)$, and put
    $R=\{\iota(i+1,j+1):0\leq i<W,\ 0\leq j<H\}$.
    Every two labelled unit-step paths in $D$ with the same lifted
    endpoints have projections related by $\sim_R$.
    This is the rectangular bulk specialization of
    \cite[Lemma~6.14]{Schuch2010PEPS}.
\end{theorem}
\begin{proof}
    Let $C_{i,j}$ run along the bottom row from $(0,0)$ to $(i,0)$,
    then up to $(i,j)$. Interchanging an east edge with a column of
    height $j$ uses exactly the $j$ adjacent squares. Hence
    $C_a e\sim_R C_b$ for an east edge $e:a\to b$.
    The north-edge identity holds directly; west and south follow
    by cancelling an edge with its reverse. Induction on the steps
    of $p:a\to b$ gives $C_a p\sim_R C_b$.
    Thus $C_a p\sim_R C_a q$; prepend $C_a^{-1}$ and cancel
    its successive backtracks. The bounds $W<w$ and $H<h$ make
    $\iota$ injective on $D$. No periodic-endpoint identification
    replaces equality of the lifted endpoints.
\end{proof}

% Ink-to-index: all dotted arrows are dual geometric routes, with no tensor
% indices or physical legs. In each panel a=(0,0), b=(1,2); vertical drawing
% rows decrease as lifted j increases. The two square centers s_1=(1,1),
% s_2=(1,2) are primal-site labels in the lower-left-plaquette convention.
% The rectangle is two squares high and one square wide. The left route is
% N,N,E; the right route E,N,N. Only these two primal sites are swept.
% This depicts the column/east exchange, not the whole arbitrary-path proof.
\begin{tenkzequation}
    \tenkzkernel
    \begin{tenkz}[size=m, rows={wire,wire,wire,wire,wire,wire,wire}, cols=7, bonds=none]
        \tnwire[kind=string, stroke=dotted, route=orth, dir=to,
            name=rectColumnEast, via={(2,2)}]{(6,2)}{(2,6)}
        \tnmark[form=label, label pos=225]{(6,2)}{$a$}
        \tnmark[form=label, label pos=45]{(2,6)}{$b$}
        \tnmark[form=label]{(5,4)}{$s_1$}
        \tnmark[form=label]{(3,4)}{$s_2$}
    \end{tenkz}
    $\quad\sim_R\quad$
    \begin{tenkz}[size=m, rows={wire,wire,wire,wire,wire,wire,wire}, cols=7, bonds=none]
        \tnwire[kind=string, stroke=dotted, route=orth, dir=to,
            name=rectEastColumn, via={(6,6)}]{(6,2)}{(2,6)}
        \tnmark[form=label, label pos=225]{(6,2)}{$a$}
        \tnmark[form=label, label pos=45]{(2,6)}{$b$}
        \tnmark[form=label]{(5,4)}{$s_1$}
        \tnmark[form=label]{(3,4)}{$s_2$}
    \end{tenkz}
\end{tenkzequation}

